% GENERATED FILE. Edit canonical lessons, then run npm run generate:books.
% canonical-source-hash: fnv1a64:c4e44649cd62f006
% canonical-lessons: LA-C01-salve, LA-W01-salve-guided-copy, LA-W01-salve-delayed-copy, LA-W01-salve-dictation, LA-C01-ave, LA-C01-vale, LA-C01-gratias, LA-C01-ita-non, LA-C01-practice

\chapter{Greetings}
\label{ch:greetings}
\hlchaptermodality{1}

\begin{chapteropening}
\textbf{By the end of this chapter:} \emph{I can greet someone in Latin, thank them, answer yes or no, and say goodbye.}

The whole doorway sequence on cue: write salvē independently from its sound, use salvē or avē on arrival, say grātiās tibi agō, answer ita or nōn, and use valē on the way out.
\end{chapteropening}

\section[salvē / salvēte]{salvē / salvēte — \textquotedblleft{}hello\textquotedblright{} (be well)}

\label{lesson:LA-C01-salve}



\begin{quote}
The first Latin word — and it already contains \emph{save}, \emph{safe}, and \emph{salute}.
\end{quote}

\begin{sounds}
Latin is read as written, every letter sounded. A \textbf{macron} (ā ē ī ō ū) marks a \emph{long} vowel — held longer, same quality. Classical \textbf{v = w}. So \emph{salvē} = \textbf{SAL-weh}, not \textquotedblleft{}sal-vay.\textquotedblright{}
\end{sounds}

\begin{cousinweb}
\emph{Salvē} (to one person) / \emph{salvēte} (to several) is the imperative of \emph{salvēre}, \textquotedblleft{}to be in good health,\textquotedblright{} from \emph{salvus} \textquotedblleft{}safe, sound, whole.\textquotedblright{} To greet in Latin is to \textbf{wish someone health}. That root fans out into English — \textbf{save, safe, salvage, salvation, salute, salubrious}, the toast \textbf{\textquotedblleft{}salut!\textquotedblright{}} — and into the Romance children as the word for health itself: Spanish \emph{salud}, Italian \emph{salute}, French \emph{salut}.
\end{cousinweb}

\begin{grammarlens}[title={Grammar Lens}]
\emph{salvē} vs. \emph{salvēte}: Latin changes the \textbf{ending} to mark one listener (\emph{-e}) versus many (\emph{-ēte}). This is the feature that defines Latin — meaning carried in endings, not in word order or extra words. English lost almost all of it; its Romance grandchildren kept some. You will meet it in every chapter.
\end{grammarlens}

\subsection*{Writing — follow one greeting you can see}
Keep \textbf{salvē} visible. First point to the line over the final \textbf{ē}: that macron says to hold the vowel longer. Trace the whole word once, from \textbf{s} through \textbf{ē}, then say \emph{SAL-weh}. Stop there.

This is not a spelling test. Your eye and hand are only following a greeting whose sound and meaning are already on the page.

\subsection*{Your turn}
Say these aloud:

\begin{itemize}
\raggedright
  \item greet one friend — \emph{salvē} (SAL-weh)
  \item greet a whole room — \emph{salvēte}
  \item two English words from \emph{salvus} (save, safe, salute, salvation…)
\end{itemize}

\begin{tcolorbox}[breakable,colback=teal!4,colframe=teal!35!black,title={Before you move on}]
What does \emph{salvē} literally wish, and how do you greet more than one person? (\textquotedblleft{}Be well / be in good health\textquotedblright{}; \emph{salvēte}, with the plural \emph{-ēte} ending.)
\end{tcolorbox}
\section[salvē / salvēte]{salvē — copy one greeting with help}

\label{lesson:LA-W01-salve-guided-copy}



\begin{quote}
Point to \textbf{salvē} and say it once: \emph{SAL-weh}. The model stays on the page.
\end{quote}

\subsection*{Writing — one guided copy}
Copy \textbf{salvē} once on the line below it. Write \textbf{salve} first, then add the short macron over the final \textbf{ē}. Look back at the model whenever you want.

Say the word as you finish. The macron belongs to the vowel: it records the long sound you already heard; it is not a new letter and not something to guess.

\begin{tcolorbox}[breakable,colback=teal!4,colframe=teal!35!black,title={Before you move on}]
Compare your copy with the model. If one mark differs, repair only that mark. No hiding, no spelling from memory, and no second word today.
\end{tcolorbox}
\section[salvē]{salvē — look, hide, write, check}

\label{lesson:LA-W01-salve-delayed-copy}



\begin{quote}
Say \textbf{salvē} once: \emph{SAL-weh}, “be well.” Its sound and meaning are already familiar. Today only the amount of visual help changes.
\end{quote}

\subsection*{Writing — delayed copy}
Look at \textbf{salvē} for five seconds. Notice the five letters and the short macron over the final \textbf{ē}. Cover the model completely. From memory, write the greeting once. Say \emph{SAL-weh} as you finish.

Uncover the model and compare one feature at a time: \textbf{s-a-l-v-ē}. If one feature differs, repair only that feature. This is one gentle memory step, not a speed test.

\begin{tcolorbox}[breakable,colback=teal!4,colframe=teal!35!black,title={Before you move on}]
Writing pass: \textbf{1/1} only when the covered-model attempt contains \textbf{s-a-l-v-ē} with the macron. Saying the greeting correctly does not replace the writing point.
\end{tcolorbox}
\section[heard greeting]{Hear the greeting; write it independently}

\label{lesson:LA-W01-salve-dictation}



\begin{quote}
Cover the previous lesson and the answer box below. Hear this cue without looking at Latin spelling: \textbf{“Greet one person: SAL-weh — literally, be well.”}
\end{quote}

\subsection*{Writing — dictation and transcription}
Write the one-word Latin greeting from the sound and meaning cue. Remember that the final vowel is long, so its written vowel carries the macron. Do not copy from an earlier page. Stop after one attempt.

\begin{tcolorbox}[breakable,colback=teal!4,colframe=teal!35!black,title={Before you move on}]
Now uncover the answer: \textbf{salvē}.

\begin{itemize}
\raggedright
  \item Writing: \textbf{1/1} for all five letters plus the final macron; otherwise \textbf{0/1}.
  \item Listening: \textbf{1/1} if the cue selected the one-person greeting; otherwise \textbf{0/1}.
\end{itemize}

Each skill must pass on its own. A correct listening choice cannot compensate for a missing or incorrect written form.
\end{tcolorbox}
\section[avē / avēte]{avē / avēte — \textquotedblleft{}hail\textquotedblright{} (a formal greeting)}

\label{lesson:LA-C01-ave}



\begin{quote}
The ceremonious greeting — the one you still hear in \emph{Ave Maria}.
\end{quote}

\begin{sounds}
\emph{Avē} = \textbf{AH-weh} (again the \emph{w}-sound for \emph{v}). Two syllables, the second long.
\end{sounds}

\begin{cousinweb}
\emph{Avē} is the more formal greeting — \textquotedblleft{}hail\textquotedblright{} — from \emph{avēre} \textquotedblleft{}to fare well, to be eager.\textquotedblright{} It survives in English essentially unchanged inside the prayer \textbf{\textquotedblleft{}Ave Maria\textquotedblright{}} (\textquotedblleft{}Hail, Mary\textquotedblright{}), and stands in Catullus's famous \emph{avē atque valē} — \textquotedblleft{}hail and farewell\textquotedblright{} — spoken at his brother's grave, pairing the word you say on arrival with the one you say on leaving (next lesson).
\end{cousinweb}

\begin{culture}
\emph{Avē} carried rank and ritual — addressed to emperors and gods (\emph{avē, Caesar}). For an everyday \textquotedblleft{}hi,\textquotedblright{} Romans reached for \emph{salvē}; \emph{avē} was the bow, not the wave.
\end{culture}

\subsection*{Your turn}
Say these aloud:

\begin{itemize}
\raggedright
  \item read it — AH-weh
  \item where English still uses it (\emph{Ave Maria}, \textquotedblleft{}Hail Mary\textquotedblright{})
  \item which is more formal, \emph{avē} or \emph{salvē}? (\emph{avē})
\end{itemize}

\begin{tcolorbox}[breakable,colback=teal!4,colframe=teal!35!black,title={Before you move on}]
What register does \emph{avē} carry, and in what famous phrase does English keep it? (Formal / ceremonial, for emperors and gods; \emph{Ave Maria}.)
\end{tcolorbox}
\section[valē / valēte]{valē / valēte — \textquotedblleft{}goodbye\textquotedblright{} (be strong)}

\label{lesson:LA-C01-vale}



\begin{quote}
Latin's goodbye wishes you \emph{strength} — and hides inside \emph{value}, \emph{valid}, and \emph{prevail}.
\end{quote}

\begin{sounds}
\emph{Valē} = \textbf{WAH-leh}. Like \emph{salvē}, a health-and-strength wish frozen into a command.
\end{sounds}

\begin{cousinweb}
\emph{Valē} (singular) / \emph{valēte} (plural) is the imperative of \emph{valēre}, \textquotedblleft{}to be strong, to be well.\textquotedblright{} So goodbye in Latin literally wishes you \textbf{strength}. The root is an English gold mine: \textbf{value, valid, valiant, valour, valence, prevail, convalesce, ambivalent} (\textquotedblleft{}strong both ways\textquotedblright{}). And the Romance children carry \emph{valēre} as their verb \textquotedblleft{}to be worth\textquotedblright{}: Spanish \emph{valer}, Italian \emph{valere}, French \emph{valoir}.
\end{cousinweb}

\begin{grammarlens}[title={Grammar Lens}]
\emph{salvē}, \emph{avē}, \emph{valē} are all the same shape — a bare \textbf{imperative}, a command form of a verb about health or strength. Latin doesn't say \textquotedblleft{}hello/goodbye\textquotedblright{}; it \emph{orders} you to thrive. Feel the \emph{-ē} / \emph{-ēte} pair recur: know it once, and you can address one friend or a whole room in any of the three.
\end{grammarlens}

\subsection*{Your turn}
Say these aloud:

\begin{itemize}
\raggedright
  \item say goodbye to one person — \emph{valē} (WAH-leh)
  \item to several — \emph{valēte}
  \item the whole \textquotedblleft{}hail and farewell\textquotedblright{} — \emph{avē atque valē}
  \item two English words from \emph{valēre} (value, valid, valiant, prevail…)
\end{itemize}

\begin{tcolorbox}[breakable,colback=teal!4,colframe=teal!35!black,title={Before you move on}]
What does \emph{valē} literally wish, and what everyday English noun comes straight from its root? (\textquotedblleft{}Be strong / be well\textquotedblright{}; \emph{value} — from \emph{valēre} \textquotedblleft{}to be worth.\textquotedblright{})
\end{tcolorbox}
\section[grātiās (tibi) agō]{grātiās (tibi) agō — \textquotedblleft{}thank you\textquotedblright{} (I give thanks)}

\label{lesson:LA-C01-gratias}



\begin{quote}
Say this and you have said \emph{grace}, \emph{gratitude}, \emph{gracias}, and \emph{grazie} all at once.
\end{quote}

\begin{sounds}
\emph{Grātiās} = \textbf{GRAH-tee-ās}; \emph{agō} = \textbf{AH-go}. In classical Latin \emph{c} and \emph{t} before \emph{i} stay hard — \textbf{GRA-tee}, never \textquotedblleft{}GRA-shee.\textquotedblright{}
\end{sounds}

\begin{cousinweb}
The full phrase is \emph{grātiās tibi agō} — literally \textquotedblleft{}\textbf{I give thanks to you}\textquotedblright{} (\emph{grātiās} \textquotedblleft{}thanks,\textquotedblright{} \emph{tibi} \textquotedblleft{}to you,\textquotedblright{} \emph{agō} \textquotedblleft{}I do/drive\textquotedblright{}). The noun \emph{grātia} means \textquotedblleft{}favour, grace, gratitude,\textquotedblright{} from \emph{grātus} \textquotedblleft{}pleasing, thankful.\textquotedblright{} Into English: \textbf{grace, gratitude, gratis, grateful, gracious, ingratiate, congratulate}. And straight into the Romance \textquotedblleft{}thank you\textquotedblright{} you already learned — Italian \emph{grazie}, Spanish \emph{gracias}, Portuguese \emph{graças} — all just the plural \emph{grātiās}, \textquotedblleft{}thanks,\textquotedblright{} worn smooth.
\end{cousinweb}

\begin{grammarlens}[title={Grammar Lens}]
\emph{agō} means \textquotedblleft{}I do\textquotedblright{} — the \emph{-ō} ending \textbf{is} \textquotedblleft{}I,\textquotedblright{} so no separate pronoun is needed. This is Latin's great habit: the verb ending tells you \emph{who} acts. \emph{Agō} itself is a root you know: \textbf{agent, act, agenda} (\textquotedblleft{}things to be done\textquotedblright{}), \textbf{agile, navigate}.
\end{grammarlens}

\subsection*{Your turn}
Say these aloud:

\begin{itemize}
\raggedright
  \item the phrase — \emph{grātiās tibi agō}
  \item the Romance descendants of \emph{grātiās} (grazie, gracias, graças)
  \item what the \emph{-ō} on \emph{agō} already means (the pronoun \textquotedblleft{}I\textquotedblright{})
\end{itemize}

\begin{tcolorbox}[breakable,colback=teal!4,colframe=teal!35!black,title={Before you move on}]
\emph{Grātiās agō} literally says what, and why is there no word for \textquotedblleft{}I\textquotedblright{}? (\textquotedblleft{}I give thanks\textquotedblright{}; the \emph{-ō} verb-ending already carries \textquotedblleft{}I.\textquotedblright{})
\end{tcolorbox}
\section[ita / nōn]{ita / nōn — \textquotedblleft{}yes\textquotedblright{} and \textquotedblleft{}no\textquotedblright{}}

\label{lesson:LA-C01-ita-non}



\begin{quote}
A surprise: Latin had no plain word for \textquotedblleft{}yes.\textquotedblright{} Here is what Romans said instead.
\end{quote}

\begin{sounds}
\emph{Ita} = \textbf{IH-tah}; \emph{nōn} = \textbf{nōne}, with a long ō.
\end{sounds}

\begin{cousinweb}
Classical Latin had \textbf{no simple \textquotedblleft{}yes.\textquotedblright{}} To agree you said \emph{ita} (\textquotedblleft{}thus, so\textquotedblright{}), \emph{sīc} (\textquotedblleft{}so, in this way\textquotedblright{}), \emph{certē} (\textquotedblleft{}certainly\textquotedblright{}), or you just \textbf{repeated the verb} — asked \textquotedblleft{}are you coming?\textquotedblright{}, you answered \textquotedblleft{}I'm coming.\textquotedblright{} It was \emph{sīc} that hardened into the Romance \textquotedblleft{}yes\textquotedblright{}: Spanish and Italian \textbf{sí / sì}, French \emph{si} (the emphatic yes). \textquotedblleft{}No\textquotedblright{} is easier: \emph{nōn} (\textquotedblleft{}not\textquotedblright{}), from older \emph{ne-oenum} \textquotedblleft{}not one thing\textquotedblright{} — the ancient negative \textbf{ne-} that also gives English \textbf{no, not, none, never}, and which you will meet again as Sanskrit \emph{na} and Hindi \emph{nahīṇ}.
\end{cousinweb}

\begin{culture}
A language that carries so much in its verb endings barely needs a \textquotedblleft{}yes\textquotedblright{} word — the verb \emph{is} the answer. This is a real window into how Latin thinks: agreement lives in the sentence, not in a little particle bolted on.
\end{culture}

\subsection*{Your turn}
Say these aloud:

\begin{itemize}
\raggedright
  \item three Latin ways to say \textquotedblleft{}yes\textquotedblright{} (\emph{ita}, \emph{sīc}, \emph{certē})
  \item which one became Spanish/Italian \emph{sí/sì}? (\emph{sīc})
  \item the English cousins of \emph{nōn} (no, not, none, never)
\end{itemize}

\begin{tcolorbox}[breakable,colback=teal!4,colframe=teal!35!black,title={Before you move on}]
How did Romans say \textquotedblleft{}yes,\textquotedblright{} and what Romance word grew out of \emph{sīc}? (With \emph{ita}/\emph{sīc}/\emph{certē} or by repeating the verb; \emph{sí} / \emph{sì}.)
\end{tcolorbox}
\section[Practice]{Chapter 1 — Practice and Recall}

\label{lesson:LA-C01-practice}



\begin{quote}
No new words. Just the greetings — and the English and Romance words they secretly contain.
\end{quote}

\subsection*{Your turn: Rapid recall}
Say these aloud:

\begin{itemize}
\raggedright
  \item hello, one person — \emph{salvē}
  \item hail, formally — \emph{avē}
  \item goodbye, one person — \emph{valē}
  \item thank you — \emph{grātiās tibi agō}
  \item \textquotedblleft{}thus\textquotedblright{} (a way to say yes) and \textquotedblleft{}not\textquotedblright{} — \emph{ita} / \emph{nōn}
\end{itemize}

\begin{cousinweb}
Each greeting is an etymological hub. Trace them:

Say these aloud:

\begin{itemize}
\raggedright
  \item which greeting root gives \emph{save, safe, salute}? (\emph{salvus}, from \emph{salvē})
  \item which gives \emph{value, valid, prevail}? (\emph{valēre}, from \emph{valē})
  \item which gives \emph{grace, gratitude, grazie, gracias}? (\emph{grātia}, from \emph{grātiās})
\end{itemize}
\end{cousinweb}

\begin{tcolorbox}[breakable,colback=teal!4,colframe=teal!35!black,title={Before you move on}]
\begin{itemize}
\raggedright
  \item \emph{salvus} \textquotedblleft{}safe\textquotedblright{} $\to$ \textbf{save, safe, salvage, salvation, salute}
  \item \emph{valēre} \textquotedblleft{}be strong\textquotedblright{} $\to$ \textbf{value, valid, valiant, prevail, convalesce}
  \item \emph{grātia} \textquotedblleft{}grace\textquotedblright{} $\to$ \textbf{grace, gratitude, gratis, congratulate}; Romance \emph{grazie/gracias}
  \item \emph{nōn} $\leftarrow$ \emph{ne-} $\to$ English \textbf{no, not, none}; cousin of Sanskrit \emph{na}
\end{itemize}
\end{tcolorbox}

\begin{tcolorbox}[breakable,colback=teal!4,colframe=teal!35!black,title={Before you move on}]
Catullus wrote \emph{avē atque valē} at his brother's grave. Translate it, and say what each verb literally wishes. (\textquotedblleft{}Hail and farewell\textquotedblright{}; \emph{avē} \textquotedblleft{}fare well\textquotedblright{} on meeting, \emph{valē} \textquotedblleft{}be strong\textquotedblright{} on parting.)
\end{tcolorbox}
